%! Linear Functions: Slope, Forms & Graphing — Unit 1, Lesson 1.1 — Group Activity (Answer Key)
% Mirrors activity/main.tex; each \writelines{2} is answered with exactly two \ansline's.
\documentclass[10pt]{article}
\usepackage{algebra2-article}
\usepackage{algebra2-key}   % pulls in -boxes; do NOT also load -boxes
% Gym cost graph: climbs n on 0..8, cost C on 0..30 (ticks every $6).
\newcommand{\gymgrid}{%
  \draw[xstep=1, ystep=3, linegray!40, very thin] (0,0) grid (8,30);
  \draw[->, semithick, charcoal] (0,0)--(8.7,0) node[right, font=\scriptsize]{$n$};
  \draw[->, semithick, charcoal] (0,0)--(0,32) node[above, font=\scriptsize]{$C$ (\$)};
  \foreach \x in {2,4,6,8} \node[charcoal, font=\tiny, below] at (\x,0){\x};
  \foreach \y in {6,12,18,24,30} \node[charcoal, font=\tiny, left] at (0,\y){\y};}

\begin{document}

\pageheader{Unit 1, Lesson 1.1}{Group Activity --- Answer Key}

\vspace{0.08in}

\begin{headlinebox}{forestbg}
\small\textbf{Two Receipts.} Two climbing gyms charge the same way --- something to walk in the door,
then a price per climb --- but they tell you about it very differently. Work with your group. Use the
three forms from the notes, and for every number you write down, be ready to show \emph{where it came
from}.
\end{headlinebox}

\vspace{0.06in}

\begin{scenariobox}[1. Summit Gym --- They Show You the Graph]{forest}
\begin{minipage}[t]{0.56\linewidth}
\vspace{0pt}
\begin{enumerate}[label=\textbf{\alph*.}, leftmargin=1.8em, itemsep=5pt]
  \item What does a visit cost if you climb \emph{zero} times? \$\ans{$6$} \;
        What is that charge, in the story? \ans{the walk-in fee}
  \item Each extra climb adds \$\ans{$3$}. Mark on the graph where you see that, and say
        which two points you used.
        \ansline{From $(0,6)$ to $(4,18)$: \$12 over 4 climbs, so \$3 each.}
        \ansline{On the graph, one step right is one tick ($\$3$) up.}
  \item Write Summit's rule in \textbf{slope-intercept} form. \; $C(n) = $ \ans{$3n+6$}
\end{enumerate}
\end{minipage}\hfill
\begin{minipage}[t]{0.40\linewidth}
\vspace{0pt}\centering
\begin{tikzpicture}[x=0.50cm, y=0.12cm]
  \gymgrid
  \draw[very thick, forest] (0,6)--(8,30);
  \fill[charcoal] (0,6) circle (2.6pt);
  \fill[charcoal] (4,18) circle (2.6pt) node[right, font=\tiny, charcoal]{$(4,18)$};
\end{tikzpicture}\\[1pt]
{\footnotesize Summit Gym: cost of $n$ climbs}
\end{minipage}

\begin{enumerate}[label=\textbf{\alph*.}, start=4, leftmargin=1.8em, itemsep=5pt]
  \item Use your rule to price 7 climbs: \$\ans{$27$}. Check it against the graph --- do the
        rule and the picture agree? \ans{yes}
\end{enumerate}
\end{scenariobox}

\boxguard[24]
\begin{scenariobox}[2. Basecamp Gym --- All You Get Is Two Receipts]{navy}
No graph, no posted prices. A friend hands you two receipts: \textbf{3 climbs cost \$24}, and
\textbf{7 climbs cost \$32}.

\begin{enumerate}[label=\textbf{\alph*.}, leftmargin=1.8em, itemsep=5pt]
  \item Going from the first receipt to the second is \ans{$4$} more climbs for
        \$\ans{$8$} more, so the slope is \$\ans{$2$} per climb.
  \item Work back to what Basecamp charges for \emph{zero} climbs.
    \begin{work}
      \text{3 climbs} &= \$24 \\
      \text{3 climbs alone cost } 3(\$2) &= \$6 \\
      \text{so the walk-in fee is } \$24 - \$6 &= \$18
    \end{work}
    Walk-in fee: \$\ans{$18$} \; Rule: \; $D(n) = $ \ans{$2n+18$}
  \item \textbf{Priya never found the walk-in fee.} She started from the receipt she had and wrote
        $$D - 24 = 2(n - 3).$$
        Which \emph{form} is that? \ans{point-slope} \; Is her rule the same line as yours, or a
        different one? Test both at $n = 7$, then tidy hers up.
    \begin{work}
      D - 24 &= 2(n-3) \\
      D - 24 &= 2n - 6 \\
           D &= 2n + 18
    \end{work}
    \ansline{The same line. At $n=7$ both give \$32, and expanding hers gives $D=2n+18$ exactly.}
    \ansline{She used the point she had; we used the intercept we had to work out. Same rule.}
  \item Now put Basecamp in \textbf{standard} form ($Ax+By=C$, integers, $A\ge0$).
        \ans{$2n - D = -18$}
\end{enumerate}
\end{scenariobox}

\boxguard[22]
\begin{scenariobox}[3. Which Gym Is Cheaper?]{forest}
Summit: $C(n)=3n+6$. \quad Basecamp: $D(n)=2n+18$.
\begin{enumerate}[label=\textbf{\alph*.}, leftmargin=1.8em, itemsep=5pt]
  \item Cheaper for \textbf{2 climbs}? \ans{Summit, \$12} \; For \textbf{15 climbs}?
        \ans{Basecamp, \$48}
  \item Find the number of climbs where the two cost the same.
    \begin{work}
      3n + 6 &= 2n + 18 \\
           n &= 12 \\
      C(12) = D(12) &= \$42
    \end{work}
    So the answer switches at $n=\ans{$12$}$ climbs, and both cost \$\ans{$42$} there.
  \item On the same axes, Basecamp crosses the vertical axis \ans{higher} than Summit and is
        \ans{less} steep. Are the two lines parallel? \ans{no} \; How do you
        know from the slopes alone?
        \ansline{Parallel lines have equal slopes; these are $3$ and $2$, so the lines are not parallel.}
        \ansline{Different slopes means they must cross exactly once --- at $n=12$.}
\end{enumerate}
\end{scenariobox}

\boxguard[24]
\begin{scenariobox}[4. Model It --- The Drone Descent]{navy}
A drone comes down at a steady rate. At $t = 2$ seconds it is $46$~m up; at $t = 6$ seconds it is
$30$~m up. Let $h$ be its height after $t$ seconds.
\begin{enumerate}[label=\textbf{\alph*.}, leftmargin=1.8em, itemsep=5pt]
  \item Find the rate.
    \begin{work}
      m &= \frac{30 - 46}{6 - 2} \\
        &= \frac{-16}{4} \\
        &= -4
    \end{work}
    So the drone falls \ans{$4$ m per second}, and the sign is negative because \ans{the height is falling}
  \item Write the line in \textbf{point-slope} form using $(2,46)$: \ans{$h-46=-4(t-2)$}
        \\[2pt] Now put it in \textbf{slope-intercept} form.
    \begin{work}
      h - 46 &= -4(t - 2) \\
      h - 46 &= -4t + 8 \\
           h &= -4t + 54
    \end{work}
    The $54$ means \ans{the height at $t=0$, where it began} --- a moment nobody measured. \\[2pt]
    In \textbf{standard} form: \ans{$4t + h = 54$}
  \item Suppose the drone had begun the same descent from $60$~m instead. In $h = mt + b$, which
        number changes and which stays put --- and what does that do to the graph?
        \ansline{$b$ changes from $54$ to $60$; the rate $m=-4$ stays put.}
        \ansline{The line shifts \emph{up} $6$~m --- same steepness, higher start.}
\end{enumerate}
\end{scenariobox}

\end{document}
